\chapter{Results from Bronze, sapphire, and Nb substrates}

    \section{Nb$_3$Sn grown via a Nb deposition on bronze}
    
    This work describes a rapid Nb$_3$Sn film growth when Nb is deposited on hot bronze, producing an immediate conversion to Nb$_3$Sn. This reaction is compared to a post-reacted Nb$_3$Sn film on bronze. This rapid reaction is not available to solid-state reactions. A factor of 12 faster growth rate for $\approx$ 700 nm thick films was observed, in comparison to deposition of Nb onto bronze at 200 $^{\circ}$C (where no reaction occurs) followed by a post-reaction at 700 $^{\circ}$. The hot bronze Nb$_3$Sn films had a columnar grain structure similar to Thornton structure zone 2 shown in Figure \ref{fig:Thornton}, whereas the post-reaction method produced smaller equiaxed grains characteristic of reactions in wires as shown in Figure \ref{fig:HotBronzeonBronze}. The films made in this work are given in Table \ref{tbl:HotBronzeonBronze}. 

    \begin{figure}
        \centering
        \includegraphics[width=.7\textwidth]{images/Bronze/Thornton.png}
        \caption{Possible microstructure available for sputtering depositions. The thin films made with the rapid hot-bronze technique resembled the columnar structure shown in zone 2, and the films made with the post-reaction method resemble zone 3 \cite{andersStructureZoneDiagram2010}}
        \label{fig:Thornton}
    \end{figure}

    Samples showed a smooth Nb$_3$Sn surface with no cracks and high tin content. The critical temperature was consistent with tin content and strain, significantly reducing to 14 K – 16 K from the 18.3 K of bulk stoichiometric Nb$_3$Sn. This depression in $T_c$ to 16 K is consistent with strain-induced CTE mismatch between Nb$_3$Sn and bronze substrate. A depression further to 14 K from 16 K is due to sub 25\% at. Sn in the Nb$_3$Sn The hot-bronze approach can be scaled to applications such as SRF cavities and devices as it produces smooth, dense, crack-free, and stoichiometric Nb$_3$Sn coatings with a relatively easy process \cite{withanageRapidNbSn2021}. 

    Understanding Nb growth on bronze is crucial, as it offers a simple starting point for studying Nb$_3$Sn formation. We observed significant differences in the microstructure and electromagnetic properties of Nb$_3$Sn formed through various methods: sputtering Nb onto hot and cold bronze substrates, with and without a post-reaction diffusion. However, bronze's thermal conductivity is lower than Cu's. Therefore, we aim to leverage these insights from the bronze system to facilitate a similar Nb-bronze reaction on a Cu substrate, ultimately achieving the desired properties with improved thermal management.
    
    \begin{table}
        \caption{Summary of sample preparation details and properties. \cite{withanageRapidNbSn2021}}
        \includegraphics[width=\textwidth]{images/Bronze/Bronzesamples.png}
        \label{tbl:HotBronzeonBronze}
    \end{table}
    
    
    \begin{figure}
        \centering
        \includegraphics[width=.5\textwidth]{images/Bronze/HotBronzeonBronze.png}
        \caption{ SEM-SE images of cross-sections of Nb$_3$Sn films made at 715 $^{\circ}$C by deposition of Nb on hot bronze (a) and via post-reaction (b) \cite{withanageRapidNbSn2021}.}
        \label{fig:HotBronzeonBronze}
    \end{figure}
    

    \begin{figure}
        \centering
        \includegraphics[width=.5\textwidth]{images/Bronze/BronzeReactionsTc.png}
        \caption{ a) T$_c$ measurements of bronze substrate coupon samples made with a 715C 6 hr post-reaction, a 715C hot bronze reaction, and both a hot bronze and post reaction. b) The onset of the T$_c$ transitions. }
        \label{fig:BronzeReactionsTc}
    \end{figure}
    
    

    %Show Tc curves??

\section{Nb$_3$Sn grown via Nb deposited on a bronze layer on a sapphire substrate}
        
        To eliminate the impact diffusion into the substrate had on Nb$_3$Sn growth, a Cu-Sn 2 \si{\um} layer of low Sn 13\% wt. was deposited onto an inert sapphire substrate via thermal evaporation. Nb was sputtered onto this structure while at 700$^{\circ}$C. This approach allowed us to study the hot bronze-Nb reaction within a thin micron-scale layer without chemical interference from the substrate. With Sn volume held constant with a passive substrate, we studied the strain on the Nb$_3$Sn film from the interface with the substrate and Cu-Sn film. We measured the critical temperature of the film still confined to the sapphire, peeled from the sapphire, and peeled off the substrate and etched from the bronze, leaving a free-standing Nb$_3$Sn film. Figure \ref{fig:combinedTc} a) shows the full $T_c$ curves and b) a cross-section of the film still confined to the sapphire as described below:
        
        \begin{enumerate}
            \item \textbf{On sapphire:} The film exhibited a narrow transition width with a $T_c$ onset of 17.1 K and a small secondary transition at 8 K, indicating the presence of unreacted Nb.
            \item \textbf{Peeled from sapphire:} $T_c$ onset decreased slightly to 16.8 K, and the Nb signal at 8 K diminished further.
            \item \textbf{Peeled and etched:} After removing residual Cu-Sn with APS, the final $T_c$ onset increased to 17.6 K, but the transition width broadened, and the Nb signal at 9 K grew dramatically because it was no longer screened by the Nb$_3$Sn signal. This suggests that Nb$_3$Sn strain decreased, but purity after etching decreased.
        \end{enumerate}
        
        \begin{figure}
            \centering
            \begin{subfigure}{0.45\textwidth}
            \centering
            \includegraphics[width = \textwidth]{Peeled.png}
            \caption{Magnetization curve}
            \label{fig:Peeled}
            \end{subfigure}
            \begin{subfigure}{0.45\textwidth}
            \centering
            \includegraphics[width = \textwidth]{crossSSub13CuSnNb3Sn.png}
            \caption{SEM/FIB cross section}
            \label{fig:crossSSub13CuSnNb3Sn}
            \end{subfigure}
            \caption[Nb$_3$Sn on sapphire experiment, to test the effect of strain from substrate and Cu-Sn layer.]{a) Compares $T_c$ of a sample confined to a sapphire substrate, peeled away, and then etched with APS after peeling. The cross-section of this confined sample is shown in b). This thin film Nb$_3$Sn 300nm
            %(check??)
            on 2 µm of 13\% Cu-Sn on Sapphire is reacted with the hot bronze method. $\mathrm{Nb_{3}Sn}$ on sapphire has been made to study the substrate materials effects on the properties of the film. }
            \label{fig:combinedTc}
        \end{figure}
        

        %We can see void space between the film and the sapphire substrate, which might reduce the strain in the film. 
        
        The CTE of sapphire is smaller than Nb$_3$Sn, which will reduce its impact on depressing the $T_c$. The critical temperature of 17.1 K for the film still on sapphire shows that the film had less strain than the films on bronze. When we peeled the sample off the sapphire substrate, the critical temperature decreased to 16.8 K. The Cu-Sn near the Nb$_3$Sn was now dominating the strain, and since Cu-Sn has a large CTE, the strain was increased, lowering $T_c$. When the film was etched, the Cu-Sn film was removed, and the Nb$_3$Sn film was no longer confined to the Cu-Sn film, thereby reducing strain and increasing $T_c$ to 17.6 K. The strain dependence on $T_c$ for Nb$_3$Sn is given in Figure \ref{fig:TcvsStrain}. The $T_c$ transition width also increased, which means some of the good quality Nb$_3$Sn was lost from the etching. The Nb signal at 8 K was also much stronger for the peeled and etched film, which might be explained by the etching opening up areas of pure Nb that the Nb$_3$Sn screened before etching. This experiment helps us understand the impact of CTE on the strain state of the Nb$_3$Sn film, both coming from the thin films in contact with the superconducting film and the substrate. One can refer to the $T_c$ vs strain plot in Figure \ref{fig:TcvsStrain}.
    \clearpage
